\chapter{Experimental Platform and Mathematical Modelling}
\label{ch:platform_modelling}
\label{ch:modelling}

\input{chapters/sections/platform_hardware}
\input{chapters/sections/notation_signals}

\section{Modelling Assumptions}
The controller-design model assumes rigid links, frictionless joints, an
abstract commanded base acceleration $\ddot{\theta}$, no missed steps or
stalls, no sensor glitches, and no actuator saturation. These assumptions
isolate the dominant pendulum dynamics. They do not describe the complete
experimental plant, so all measured results are interpreted together with the
implementation constraints in Chapter~\ref{ch:controller_implementation}.

\section{Nonlinear Dynamic Model}
Using the Euler--Lagrange method with $\theta$ as the arm angle and $\alpha=0$
at upright gives the coupled model
\begin{align}
(J_0+J_1\sin^2\alpha)\ddot{\theta}
{}+K\cos\alpha\,\ddot{\alpha}
{}+J_1\sin(2\alpha)\dot{\theta}\dot{\alpha}
-K\sin\alpha\,\dot{\alpha}^2 &= \tau, \\
K\cos\alpha\,\ddot{\theta}
{}+J_1\ddot{\alpha}
-\frac{1}{2}J_1\sin(2\alpha)\dot{\theta}^2
-G\sin\alpha &= 0.
\end{align}
The structure is consistent with established Furuta-pendulum derivations
\cite{Cazzolato2011}. The complete coordinate construction, energy
expressions, and Euler--Lagrange algebra are retained in
Appendix~\ref{app:full_derivation}.

Solving the pendulum equation for the acceleration-input representation used
by the nonlinear controllers gives
\begin{equation}
\ddot{\alpha}
=A\sin\alpha
+\sin\alpha\cos\alpha\,\dot{\theta}^{2}
-B\cos\alpha\,\ddot{\theta}.
\label{eq:nonlinear_reduced}
\end{equation}

\section{Upright Linearization and Numerical Constants}
For $|\alpha|\ll 1$, use $\sin\alpha\approx\alpha$,
$\cos\alpha\approx1$, and neglect products of small quantities. The reduced
upright relation becomes
\begin{equation}
\ddot{\alpha}=A\alpha-B\ddot{\theta}.
\label{eq:upright_reduced}
\end{equation}

The measured geometry and mass distribution give the nominal constants in
Table~\ref{tab:model_constants}.
\begin{table}[htbp]
  \centering
  \caption{Nominal model parameters and constants used for controller design.}
  \label{tab:model_constants}
  \begin{tabular}{lll}
    \toprule
    Quantity & Value & Meaning \\
    \midrule
    $J_0$ & \SI{1.104e-3}{\kilogram\metre\squared} & effective arm-side inertia \\
    $J_1$ & \SI{1.021e-4}{\kilogram\metre\squared} & pendulum pivot inertia \\
    $K$ & \SI{1.993e-4}{\kilogram\metre\squared} & inertial coupling constant \\
    $G$ & \SI{1.029e-2}{\newton\metre} & gravity constant \\
    $A=G/J_1$ & \SI{100.8}{\per\second\squared} & unstable upright coefficient \\
    $B=K/J_1$ & 1.952 & base-acceleration coupling \\
    $k_s$ & \SI{4.444}{\step\per\degree} & angle-to-step conversion \\
    \bottomrule
  \end{tabular}
\end{table}

Thus,
\[
\ddot{\alpha}=100.8\,\alpha-1.952\,\ddot{\theta}.
\]
These are model constants, not evidence that the stepper reproduces the
commanded acceleration exactly.

\section{Model Use and Limitations}
Equation~\eqref{eq:upright_reduced} is used for linear state-feedback design,
while Equation~\eqref{eq:nonlinear_reduced} is used for the two sliding-mode
controllers. The model is not independently validated through a parameter
identification experiment in this thesis. Its role is to provide a coherent
design model; the experimental chapter evaluates the complete closed-loop
system rather than claiming plant-model equivalence.
